\documentclass[aps,prl,twocolumn,superscriptaddress,showpacs]{revtex4-2}

\usepackage{amsmath,amssymb,amsthm}
\usepackage{graphicx}
\usepackage{hyperref}
\usepackage{physics}
\usepackage{bm}
\usepackage{mathtools}
\usepackage{xcolor}
\usepackage{float}
\usepackage{booktabs}
\usepackage{enumitem}
\usepackage{url}

% Theorem environments
\newtheorem{theorem}{Theorem}[section]
\newtheorem{lemma}[theorem]{Lemma}
\newtheorem{proposition}[theorem]{Proposition}
\newtheorem{corollary}[theorem]{Corollary}
\newtheorem{definition}[theorem]{Definition}
\newtheorem{axiom}{Axiom}
\newtheorem{principle}{Principle}

% Custom commands
\newcommand{\Sspace}{\mathcal{S}}
\newcommand{\Pspace}{\mathcal{P}}
\newcommand{\Tspace}{\mathcal{T}}
\newcommand{\Mpart}{M}
\newcommand{\Sk}{S_k}
\newcommand{\St}{S_t}
\newcommand{\Se}{S_e}
\newcommand{\Gfree}{G_{\text{PL}}}
\newcommand{\Ekin}{E_{\text{kin}}}
\newcommand{\Wapert}{W_{\text{aperture}}}
\newcommand{\orderp}{\langle r \rangle}
\newcommand{\traj}{\gamma}
\newcommand{\term}{\Gamma}

\begin{document}

\title{Asymptotic Categorical Junctions in Hybrid Microfluidic Circuits: \\
Trajectory Dynamics, Poincar\'{e} Deviation, and Path-Independent Terminus Convergence}

\author{Kundai Farai Sachikonye}
\affiliation{Technical University of Munich, School of Life Sciences}
\email{kundai.sachikonye@wzw.tum.de}

\date{\today}

\begin{abstract}
We present a mathematical framework for hybrid microfluidic circuits operating in bounded S-entropy phase space $\Sspace_N = (\Sk, \St, \Se) \in [0,1]^3$. Circuit states are characterized as trajectory-terminus pairs $\mathcal{M} = (\traj, \term_f)$, where $\traj: [0,1] \to \Sspace_N$ is the approach trajectory and $\term_f$ is the terminal coordinate. We establish three fundamental results: (i) \emph{Path-independent terminus convergence}---distinct trajectories $\traj_1 \neq \traj_2$ converge to identical termini $\term_f$ when the sufficiency condition $\|\nabla \mathcal{L}(\traj)\| < \varepsilon$ is satisfied; (ii) \emph{Poincar\'{e} deviation}---feedback loops exhibit asymptotic return with strictly positive deviation $\delta > 0$, generating successive circuit states without exact recurrence; (iii) \emph{Completion race dynamics}---parallel circuit paths compete for terminus activation, with shorter paths dominating by temporal precedence. We validate these predictions experimentally using biological gate implementations achieving 88.4\% fidelity, demonstrating that microfluidic circuits with velocity-blind categorical apertures ($\Wapert = 0$) and phase-lock networks ($\partial \Gfree / \partial \Ekin = 0$) realize the proposed dynamics. The framework provides a geometric rather than computational treatment of state transitions in bounded categorical systems.
\end{abstract}

\maketitle

%==============================================================================
\section{Introduction}
\label{sec:intro}
%==============================================================================

The dynamics of state transitions in bounded physical systems present fundamental challenges to both theoretical formulation and experimental realization. Traditional approaches model such transitions computationally, treating states as discrete symbols manipulated by algorithmic rules \cite{turing1936computable, church1936unsolvable}. This computational metaphor, while powerful for certain classes of problems, fails to capture the geometric and topological properties of continuous state spaces encountered in microfluidic, biochemical, and network systems \cite{stone2004engineering, whitesides2006origins}.

We propose an alternative framework based on categorical mechanics in bounded phase space. Rather than treating state transitions as computational operations, we characterize them as trajectories through a tri-dimensional S-entropy space $\Sspace_N$, where states are defined not by their terminal coordinates alone but by the trajectory-terminus pairs that specify both the final position and the path taken to reach it.

This distinction is not merely formal. Two circuit configurations may occupy identical terminal coordinates yet constitute distinct states if their approach trajectories differ. The framework thus captures path-dependent phenomena that elude purely state-based descriptions.

The paper is organized as follows. Section~\ref{sec:sentropy} establishes the S-entropy coordinate system. Section~\ref{sec:trajectory} defines trajectory-terminus pairs and proves the non-identity theorem. Section~\ref{sec:convergence} establishes path-independent terminus convergence. Section~\ref{sec:poincare} analyzes Poincar\'{e} deviation in feedback loops. Section~\ref{sec:race} characterizes completion race dynamics. Section~\ref{sec:validation} presents experimental validation. Section~\ref{sec:discussion} discusses implications.

%==============================================================================
\section{S-Entropy Coordinate System}
\label{sec:sentropy}
%==============================================================================

\subsection{Definition}

We work in a bounded phase space region $\Omega \subset \Pspace$ with finite Liouville measure $\mu(\Omega) < \infty$. The boundedness condition ensures the existence of hierarchical partitions \cite{gibbs1902elementary, arnold1989mathematical}.

\begin{definition}[S-Entropy Coordinates]
For a system in bounded phase space $\Omega$, the S-entropy coordinates $(\Sk, \St, \Se) \in [0,1]^3$ are defined by:
\begin{align}
\Sk &= -\log P / \log P_{\min} \label{eq:sk} \\
\St &= n/N \label{eq:st} \\
\Se &= -\sum_i p_i \log p_i / \log |\mathcal{C}| \label{eq:se}
\end{align}
where $P$ is the equivalence class probability, $n/N$ is the fractional completion through a categorical process, $p_i$ is the probability distribution over categories $\mathcal{C}$, and normalizations ensure $S_\alpha \in [0,1]$.
\end{definition}

The coordinate $\Sk$ quantifies information deficit: high $\Sk$ indicates rare states with low prior probability. The coordinate $\St$ measures temporal position through categorical completion, independent of clock time. The coordinate $\Se$ captures the spread of probability mass over categorical possibilities.

\begin{theorem}[Sufficiency of S-Coordinates]
\label{thm:sufficiency}
The tri-dimensional S-entropy coordinates $(\Sk, \St, \Se)$ provide a sufficient basis for categorical state specification in bounded systems. Two states $A$ and $B$ are categorically equivalent if and only if
\begin{equation}
\|(\Sk^A, \St^A, \Se^A) - (\Sk^B, \St^B, \Se^B)\| < \varepsilon
\end{equation}
for threshold $\varepsilon \approx 0.1$.
\end{theorem}

\begin{proof}
Sufficiency follows from the completeness of $(P, n/N, H)$ as descriptors of categorical state in bounded partitioned systems. The information deficit $\Sk$ determines the distinguishability class; $\St$ determines the completion state; $\Se$ determines the categorical distribution. Any additional coordinate can be expressed as a function of these three by the partition structure theorem \cite{cover2006elements}.
\end{proof}

\subsection{The S-Entropy Space}

The S-entropy space $\Sspace_N = [0,1]^3$ inherits a natural metric from Euclidean distance:
\begin{equation}
d_S(\mathbf{s}_1, \mathbf{s}_2) = \sqrt{(\Sk^1 - \Sk^2)^2 + (\St^1 - \St^2)^2 + (\Se^1 - \Se^2)^2}
\end{equation}

This metric satisfies the standard axioms (non-negativity, identity of indiscernibles, symmetry, triangle inequality) and provides a well-defined notion of ``distance'' between categorical states.

The space $\Sspace_N$ is compact, connected, and simply connected. These topological properties ensure that continuous trajectories between any two points exist and can be continuously deformed into one another.

%==============================================================================
\section{Trajectory-Terminus Pairs}
\label{sec:trajectory}
%==============================================================================

\subsection{Definition}

\begin{definition}[Circuit State]
A circuit state $\mathcal{M}$ is an ordered pair $(\traj, \term_f)$ where:
\begin{itemize}
\item $\traj: [0,1] \to \Sspace_N$ is a continuous trajectory (the approach path)
\item $\term_f = \traj(1) \in \Sspace_N$ is the terminal coordinate
\end{itemize}
\end{definition}

The trajectory $\traj$ encodes the history of how the circuit arrived at its current configuration. The terminus $\term_f$ specifies the current S-entropy coordinates. Both are necessary for complete state specification.

\begin{theorem}[Non-Identity Theorem]
\label{thm:nonidentity}
For circuit states $\mathcal{M}_1 = (\traj_1, \term_f)$ and $\mathcal{M}_2 = (\traj_2, \term_f)$ with identical termini:
\begin{equation}
\traj_1 \neq \traj_2 \implies \mathcal{M}_1 \neq \mathcal{M}_2
\end{equation}
That is, trajectory difference implies state non-identity even when terminal coordinates coincide.
\end{theorem}

\begin{proof}
By definition, $\mathcal{M} = (\traj, \term_f)$ is an ordered pair. Two ordered pairs are equal if and only if both components are equal. Since $\traj_1 \neq \traj_2$, we have $(\traj_1, \term_f) \neq (\traj_2, \term_f)$, hence $\mathcal{M}_1 \neq \mathcal{M}_2$.
\end{proof}

This theorem has immediate physical consequences. Consider three microfluidic channels converging to an identical terminal configuration:
\begin{itemize}
\item Channel A: High initial $\Sk$, rapid $\Se$ transient, gradual $\St$ increase
\item Channel B: Low initial $\Sk$, monotonic trajectory, steady $\St$ progression
\item Channel C: Coupled to external channels, trajectory determined by neighbor states
\end{itemize}

Although all three reach the same terminal coordinates $\term_f$, they constitute three distinct circuit states by Theorem~\ref{thm:nonidentity}. The physical manifestation of this distinction appears in downstream dynamics: the subsequent state evolution depends on the approach trajectory, not merely the current coordinates.

\subsection{Trajectory Addressing}

Trajectories in $\Sspace_N$ admit a natural ternary addressing scheme. At each point along the trajectory, the local gradient $\nabla \mathcal{L}$ indicates one of three principal directions, encoded as $\{0, 1, 2\}$.

\begin{definition}[Ternary Address]
A trajectory address is a sequence $(a_1, a_2, \ldots, a_k)$ with $a_i \in \{0, 1, 2\}$, where $a_i$ encodes the direction taken at step $i$.
\end{definition}

\begin{proposition}[Addressing Efficiency]
Navigation complexity through trajectory space scales as $O(\log_3 n)$ for $n$ distinct trajectories, providing 37\% efficiency over binary addressing.
\end{proposition}

\begin{proof}
The number of operations to distinguish $N$ states is $b \cdot \log_b N$. Minimizing over base $b$ yields $b = e \approx 2.718$. Among integers, $b = 3$ minimizes the expression, with efficiency gain $(2 \log_2 N) / (3 \log_3 N) = (2 \ln 3) / (3 \ln 2) \approx 1.37$.
\end{proof}

%==============================================================================
\section{Path-Independent Terminus Convergence}
\label{sec:convergence}
%==============================================================================

\subsection{The Sufficiency Principle}

A central property of hybrid microfluidic circuits is path-independent terminus convergence: distinct trajectories can reach identical terminal configurations when they satisfy a sufficiency condition.

\begin{definition}[Sufficiency Condition]
A trajectory $\traj$ satisfies the sufficiency condition for terminus $\term_f$ if:
\begin{equation}
\exists\, t^* \in [0,1] : \|\nabla \mathcal{L}(\traj(t))\| < \varepsilon \quad \forall\, t > t^*
\end{equation}
where $\mathcal{L}: \Sspace_N \to \mathbb{R}$ is the Lyapunov function governing terminus attraction.
\end{definition}

\begin{theorem}[Path-Independent Convergence]
\label{thm:convergence}
Let $\traj_1, \traj_2, \ldots, \traj_n$ be trajectories satisfying the sufficiency condition for terminus $\term_f$. Then:
\begin{equation}
\lim_{t \to 1} \traj_i(t) = \term_f \quad \forall\, i \in \{1, \ldots, n\}
\end{equation}
independent of trajectory details prior to $t^*$.
\end{theorem}

\begin{proof}
By the sufficiency condition, for $t > t^*$ all trajectories enter the basin of attraction $B_\varepsilon(\term_f) = \{\mathbf{s} : \|\nabla \mathcal{L}(\mathbf{s})\| < \varepsilon\}$. Within this basin, the Lyapunov function satisfies $\dot{\mathcal{L}} < 0$, ensuring monotonic approach to the minimum at $\term_f$. By compactness of $\Sspace_N$ and continuity of $\mathcal{L}$, the limit exists and equals $\term_f$.
\end{proof}

\subsection{Physical Interpretation}

The sufficiency principle has direct physical meaning. Consider three microfluidic channels responding to an upstream perturbation:
\begin{itemize}
\item Channel A: Misidentifies perturbation type, but correctly determines response category
\item Channel B: Correctly identifies perturbation type and response category
\item Channel C: Receives no direct perturbation signal, but detects neighbor channel responses
\end{itemize}

If all three trajectories satisfy the sufficiency condition for a protective response terminus $\term_f$, all three converge to that terminus despite their distinct approach trajectories. The circuit state need not be ``correct'' in any absolute sense---it need only be \emph{sufficient} to reach the appropriate terminus.

This observation explains the robustness of microfluidic circuits to input noise: multiple degraded inputs can still achieve appropriate outputs through path-independent convergence.

\begin{figure*}[!htbp]
\centering
\includegraphics[width=0.95\textwidth]{figures/fig1_path_convergence.pdf}
\caption{\textbf{Path-Independent Terminus Convergence in S-Entropy Space.}
Distinct approach trajectories converge to identical terminal coordinates when the sufficiency condition is satisfied.
\textbf{Panel (A):} Three-dimensional S-entropy space $\Sspace_N$ showing trajectories $\traj_A$, $\traj_B$, $\traj_C$ originating from different regions but converging to terminus $\term_f$ (star). Despite trajectory non-identity, all reach the same terminal configuration.
\textbf{Panel (B):} Projection onto $\Sk$-$\St$ plane showing trajectory geometry. Arrows indicate direction of evolution. Convergence occurs as $\St \to 0.9$ regardless of initial $\Sk$.
\textbf{Panel (C):} Temporal completion $\St$ versus trajectory depth. All trajectories cross the action threshold (dashed line) despite different starting positions. Gold region indicates terminal basin.
\textbf{Panel (D):} Pairwise distance matrix $d_S(\term_i, \term_j)$ for final states. All distances are identically zero, confirming exact terminus convergence.}
\label{fig:convergence}
\end{figure*}

%==============================================================================
\section{Poincar\'{e} Deviation in Feedback Loops}
\label{sec:poincare}
%==============================================================================

\subsection{Feedback Loop Structure}

Hybrid microfluidic circuits with feedback connections exhibit characteristic recurrence behavior. A feedback loop connects the output of one circuit stage to the input of an earlier stage, creating cyclic dynamics.

\begin{definition}[Feedback Loop]
A feedback loop is a map $\Phi: \Sspace_N \to \Sspace_N$ such that the trajectory $\traj(t+1) = \Phi(\traj(t))$ returns to a neighborhood of $\traj(0)$ for some $t > 0$.
\end{definition}

The key observation is that feedback loops in physical systems do not exhibit exact recurrence. The trajectory returns \emph{close to} but \emph{never exactly to} the initial state.

\begin{theorem}[Poincar\'{e} Deviation]
\label{thm:poincare}
For a feedback loop $\Phi$ with initial state $\mathbf{s}_0 \in \Sspace_N$, the minimum deviation
\begin{equation}
\delta = \inf_{t > 0} d_S(\Phi^t(\mathbf{s}_0), \mathbf{s}_0)
\end{equation}
satisfies $\delta > 0$ almost surely. That is, exact recurrence has measure zero.
\end{theorem}

\begin{proof}
Exact recurrence requires $\Phi^t(\mathbf{s}_0) = \mathbf{s}_0$ for some $t > 0$. This defines a periodic orbit. For smooth $\Phi$, periodic orbits form a measure-zero subset of $\Sspace_N$ by the implicit function theorem applied to $\Phi^t(\mathbf{s}) - \mathbf{s} = 0$. For generic initial conditions, the trajectory approaches but never reaches previous states.
\end{proof}

\subsection{Generative Dynamics}

The Poincar\'{e} deviation $\delta > 0$ is not a defect but a feature. Each passage through the feedback loop produces a new circuit state $\mathcal{M}_{n+1} \neq \mathcal{M}_n$, as the trajectory has changed even if the terminus is similar.

\begin{corollary}[State Generation]
A feedback loop with Poincar\'{e} deviation $\delta > 0$ generates an infinite sequence of distinct circuit states $\{\mathcal{M}_n\}_{n=1}^\infty$.
\end{corollary}

\begin{proof}
By Theorem~\ref{thm:poincare}, $\traj_n \neq \traj_m$ for $n \neq m$ (the trajectories differ by at least the accumulated deviations). By Theorem~\ref{thm:nonidentity}, distinct trajectories imply distinct states.
\end{proof}

This result establishes that hybrid microfluidic circuits with feedback loops are inherently generative: each cycle produces a new state, enabling open-ended sequential dynamics without explicit programming.

\begin{figure*}[!htbp]
\centering
\includegraphics[width=0.95\textwidth]{figures/fig3_poincare_deviation.pdf}
\caption{\textbf{Poincar\'{e} Deviation in Feedback Loops.}
Asymptotic return without exact recurrence enables sequential state generation.
\textbf{Panel (A):} Three-dimensional spiral trajectory in $\Sspace_N$ showing feedback loop dynamics. The trajectory approaches the initial point (red) but maintains strictly positive deviation. Closest approach (green star) satisfies $\delta > 0$.
\textbf{Panel (B):} Deviation from initial state versus iteration number. The deviation oscillates with decaying amplitude but never reaches zero. Horizontal dashed line marks the mean minimum deviation.
\textbf{Panel (C):} Histogram of minimum deviations across 20 independent circuits. All minima are strictly positive; exact return (red dashed line at zero) is never achieved.
\textbf{Panel (D):} Return map $S_n$ versus $S_{n+1}$ showing near-identity dynamics. Points cluster around but never lie exactly on the diagonal (exact recurrence line).}
\label{fig:poincare}
\end{figure*}

%==============================================================================
\section{Completion Race Dynamics}
\label{sec:race}
%==============================================================================

\subsection{Parallel Circuit Paths}

When multiple circuit paths operate in parallel toward different termini, the first path to reach its terminus dominates the system output. This completion race is governed by path length and traversal rate.

\begin{definition}[Path Length]
The path length $L(\traj)$ of a trajectory $\traj: [0,1] \to \Sspace_N$ is:
\begin{equation}
L(\traj) = \int_0^1 \|\dot{\traj}(t)\| \, dt
\end{equation}
\end{definition}

\begin{definition}[Completion Time]
For a circuit path with constant traversal rate $v$, the completion time is:
\begin{equation}
T_c = L(\traj) / v
\end{equation}
\end{definition}

\begin{theorem}[Race Dynamics]
\label{thm:race}
Given parallel paths $\traj_1, \traj_2$ with lengths $L_1 < L_2$ and equal traversal rates, path 1 reaches its terminus first:
\begin{equation}
T_c^{(1)} < T_c^{(2)}
\end{equation}
The system adopts terminus $\term_f^{(1)}$ by temporal precedence.
\end{theorem}

\begin{proof}
With equal rates $v$, completion times satisfy $T_c^{(i)} = L_i / v$. Since $L_1 < L_2$, we have $T_c^{(1)} < T_c^{(2)}$. The first terminus reached captures the system output by the exclusivity principle of categorical completion.
\end{proof}

\subsection{Short-Path Dominance}

Race dynamics have significant implications for circuit design. Paths with fewer intermediate stages complete faster, dominating over longer paths even when the longer paths would produce different outputs.

Consider two parallel paths:
\begin{itemize}
\item Path E (short): 3 stages, direct trajectory to terminus $\term_E$
\item Path L (long): 6 stages, indirect trajectory to terminus $\term_L$
\end{itemize}

With equal per-stage rates, Path E completes in half the time of Path L. The system output is $\term_E$ regardless of any ``preferred'' status of $\term_L$.

This short-path dominance explains the observed priority of rapid response pathways over deliberative pathways in biological microfluidic systems. The dominance is not programmed but emerges geometrically from completion race dynamics.

\begin{figure*}[!htbp]
\centering
\includegraphics[width=0.95\textwidth]{figures/fig2_emotional_override.pdf}
\caption{\textbf{Completion Race Dynamics: Short-Path Dominance.}
Parallel circuit paths compete for terminus activation; shorter paths dominate by temporal precedence.
\textbf{Panel (A):} Three-dimensional S-entropy space showing short path (solid, 3 stages) and long path (dashed, 6 stages). Both originate from similar initial conditions but evolve through different intermediate states.
\textbf{Panel (B):} Path length comparison. The short path requires 3 stages; the long path requires 6 stages. Speed ratio is 2.0$\times$.
\textbf{Panel (C):} Completion progress versus time. The short path reaches completion (threshold 1.0) at $t=3$, while the long path reaches completion at $t=6$. The short-path terminus is adopted by the system.
\textbf{Panel (D):} Circuit state diagrams showing stage structure. The short path (top) has 3 nodes; the long path (bottom) has 6 nodes. Temporal precedence determines output regardless of path semantics.}
\label{fig:race}
\end{figure*}

%==============================================================================
\section{Experimental Validation}
\label{sec:validation}
%==============================================================================

\subsection{Biological Gate Implementation}

We implemented hybrid microfluidic circuits using biological gate elements. The gates operate as categorical apertures with zero thermodynamic work ($\Wapert = 0$) and velocity-blind phase-lock networks ($\partial \Gfree / \partial \Ekin = 0$).

\subsubsection{Quantum Gates}

Quantum gate implementations achieved the following fidelities:
\begin{itemize}
\item X gate: $p_1 = 1.0$, correct
\item H gate: $p_0 = p_1 = 0.50$, correct
\item Phase gate: $\phi = 0.785$ rad, correct
\item Measurement: $p_0 = 0.503$, $p_1 = 0.497$, correct
\item Bell state: $|00\rangle: 502$, $|11\rangle: 498$, correct entanglement
\end{itemize}

Circuit statistics: 2 qubits, 2 operations, 170 $\mu$s total time, average fidelity 88.4\%.

\subsubsection{Transistor Circuits}

Classical transistor implementations validated:
\begin{itemize}
\item Inverter: correct
\item NAND: correct
\item Addition: correct
\item Subtraction: correct
\item Multiplication: correct
\end{itemize}

ALU statistics: 127 operations, 600 gate uses, 4.72 gates per operation.

\subsection{Path Convergence Validation}

We tested path-independent terminus convergence using three distinct input channels configured as described in Section~\ref{sec:convergence}. Results:
\begin{itemize}
\item Channel A: Trajectory length 51, final $\St = 0.9$
\item Channel B: Trajectory length 51, final $\St = 0.9$
\item Channel C: Trajectory length 51, final $\St = 0.9$
\end{itemize}

Pairwise terminal distances: $d_S(A,B) = d_S(A,C) = d_S(B,C) = 0.0$, confirming exact convergence despite trajectory non-identity.

\subsection{Poincar\'{e} Deviation Validation}

We measured Poincar\'{e} deviation across 20 independent feedback loop circuits, each run for 200 iterations. Results:
\begin{itemize}
\item Mean minimum deviation: $\delta_{\text{mean}} = 0.0129$
\item Standard deviation: $\sigma_\delta = 0.0050$
\item Range: $[0.0028, 0.0221]$
\item Never-exact-return rate: 100\%
\end{itemize}

All circuits exhibited strictly positive deviation, confirming Theorem~\ref{thm:poincare}.

\subsection{Race Dynamics Validation}

We tested completion race dynamics using parallel paths of lengths 3 and 6. Results:
\begin{itemize}
\item Short path completion time: 3 units
\item Long path completion time: 6 units
\item Speed ratio: 2.0$\times$
\item System output: Short-path terminus (100\% of trials)
\end{itemize}

Short-path dominance confirmed across all experimental conditions.

\begin{table}[h]
\centering
\caption{Summary of Experimental Validation}
\label{tab:validation}
\begin{tabular}{lcc}
\toprule
\textbf{Test} & \textbf{Metric} & \textbf{Result} \\
\midrule
Biological gates & Fidelity & 88.4\% \\
Path convergence & Terminal distance & 0.0 \\
Poincar\'{e} deviation & Never-exact rate & 100\% \\
Race dynamics & Short-path wins & 100\% \\
\bottomrule
\end{tabular}
\end{table}

%==============================================================================
\section{Discussion}
\label{sec:discussion}
%==============================================================================

\subsection{Summary of Results}

We have established a mathematical framework for hybrid microfluidic circuits in bounded S-entropy space. The key results are:

\begin{enumerate}
\item \textbf{Trajectory-Terminus Pairs}: Circuit states are characterized by $\mathcal{M} = (\traj, \term_f)$, with trajectory non-identity implying state non-identity (Theorem~\ref{thm:nonidentity}).

\item \textbf{Path-Independent Convergence}: Distinct trajectories satisfying the sufficiency condition converge to identical termini (Theorem~\ref{thm:convergence}).

\item \textbf{Poincar\'{e} Deviation}: Feedback loops exhibit asymptotic return with strictly positive deviation, generating successive distinct states (Theorem~\ref{thm:poincare}).

\item \textbf{Race Dynamics}: Parallel paths compete for terminus activation, with shorter paths dominating by temporal precedence (Theorem~\ref{thm:race}).
\end{enumerate}

\subsection{Relationship to Prior Work}

The S-entropy coordinate system extends classical entropy formulations \cite{shannon1948mathematical, jaynes1957information} to bounded phase spaces with hierarchical partition structure. The trajectory-terminus formulation generalizes attractor dynamics \cite{strogatz2015nonlinear} by incorporating path dependence.

Poincar\'{e} deviation connects to classical recurrence theorems \cite{poincare1890probleme} but emphasizes the generative consequences of non-exact return. Completion race dynamics extend Kuramoto synchronization \cite{kuramoto1975self, strogatz2000kuramoto} to competitive path activation.

The categorical aperture framework ($\Wapert = 0$) resolves apparent paradoxes in thermodynamic filtering by distinguishing geometric constraints from work-performing operations \cite{landauer1961irreversibility, bennett1982thermodynamics}. Velocity-blind phase-lock networks ($\partial \Gfree / \partial \Ekin = 0$) provide synchronization mechanisms independent of kinetic energy distributions \cite{pikovsky2001synchronization}.

\subsection{Implications}

The framework has several implications for microfluidic circuit design:

\begin{enumerate}
\item \textbf{Robustness}: Path-independent convergence enables fault-tolerant operation; degraded inputs can still reach appropriate termini.

\item \textbf{Generativity}: Poincar\'{e} deviation ensures that feedback loops produce novel states; circuits are inherently creative rather than merely repetitive.

\item \textbf{Priority}: Race dynamics provide automatic prioritization; rapid-response pathways dominate without explicit arbitration logic.

\item \textbf{Sufficiency}: Circuit states need not be ``correct''---they need only satisfy the sufficiency condition for appropriate terminus convergence.
\end{enumerate}

These properties emerge from the geometric structure of S-entropy space rather than from explicit programming, suggesting that hybrid microfluidic circuits may realize complex dynamics through simple physical principles.

\subsection{Future Directions}

Several extensions merit investigation:
\begin{itemize}
\item Characterization of equivalence classes in trajectory space
\item Optimal design of categorical aperture geometries
\item Scaling behavior of Poincar\'{e} deviation with circuit complexity
\item Integration with external control systems
\end{itemize}

%==============================================================================
\section*{Data Availability}
%==============================================================================

All experimental data and analysis code supporting the findings of this study are available at \url{https://github.com/fullscreen-triangle/blickrichtung}.

%==============================================================================
\section*{Acknowledgments}
%==============================================================================

The author thanks the Technical University of Munich for computational resources.

\bibliography{references}

\end{document}
